\chapter{Guerra di rete}

«Che schifo! Non mi leccare le orecchie.»

Rocky saluta entusiasta Valentina e Alice. \`E un cane di taglia media: se Laura non gi  limasse le unghie, le sue \textit{feste} lascierebbero dei ricordi netti sulla pelle. Per fortuna sono ben arrotondate, cos\`i, tranne qualche batterio, non lascia spiacevoli graffi sulle ragazzine.\\
Laura si spoglia della tuta, un po' calda per la stagione, e si infila una magliettina.
«Non era mai successo che si allontanasse da solo. Evidentemente deve aver subito il fascino della piccola vagabonda.»

«Già... Dai, \`e andata bene...»\\

Laura guarda l’amica: «Era un sospiro quello?»
Valentina si stropiccia gli occhi e prova ad abbracciare Rocky, che si sottrae agilmente e si accuccia ai piedi del futon.

\par\medskip
«Rocky! Cos'\`e successo? State bene? Ma \`e tutto baganto!»\\
«\`E una storia lunga, ma dopo ve la raccontiamo.»\\
«Sì, ragazze, adesso abbiamo bisogno che ci aiutate. Tra trenta minuti scade la possibilità di salvare il viaggio di Caterina. Quindi abbiamo bisogno di concetrazione assoluta.»\\

Alice guarda la sorella: «E se va male? Resti a casa?»\\
«Preferiresti?» le chiede. La osserva. Alice non risponde, ma \`e Laura a parlare: «Adesso non distraiamoci. Dobbiamo completare alcune... procedure informatiche. Quindi mi raccomando.»\\
Valentina le corre vicino: «Devi hackerare un sito, tata?»

\par\medskip
Alice spalanca gli occhi, gli hacker hanno perso il diritto di cronca prima che lei nascesse, per\`o dove conta se ne parla ancora, e una ragazza intelligente non si lascia sfuggire certi dettagli: «Cosa? Sei anche un'hacker? Ma \`e fantastico!»\\
Laura non si sforza di negare, ma una precisazione \`e dovuta:
«Vista la rete in cui devo entrare diciamo che non \`e proprio hacking, \`e pi\`u archeologia digitale... Ragazze, \`e meglio se di queste cose non ne parliamo con nessuno, ok?»\\
«Vogliamo solo capire perch\'e mi hanno negato il visto» completa Caterina.\\
«Cos’\`e il visto? A proposito, ma ci stiamo dimenticando della...» A Valentina manca un grosso pezzo della storia, per\`o ha chiara in mente un'altra faccenda, peccato che il suo entusiasmo venga fermato dalla \textit{consorella minore} che la pizzica, le strizza l’occhio e si avvicina al suo orecchio: «Dai, non dire niente che facciamo una sorpresa alla fine.»\\
Valentina ricambia il gesto confidenziale: «Ok», le conferma e sigilla il patto con un pizzicotto meno trattenuto che la fa sobbalzare.

\par\medskip
«Va bene. Prepariamo gli strumenti.» Laura ha poco tempo e molto bisogno di concentrazione.\\
«Ma da quando sei una hacker?» le chiede Alice. Si sposta vicino a lei. «Questo allora \`e un sistema operativo per voi... hacker?»\\
Laura la guarda, le sorride. Il tempo stringe. «Sì, \`e un sistema dedicato...»\\
«Mi fai vedere come funziona?» Alice guarda la sorella. La contrazione delle labbra di Caterina esprime il suo rimprovero, ma poi il rimprovero diventa verbale: «Alice, abbiamo davvero poco tempo.»\\
«Ho fame, mangiamo un dolce?» Valentina torna alla carica. Nonostante la raccomandazione della consorella minore, il pensiero della crostata le stimola l’appetito.\\
Laura chiude gli occhi, prende un respiro. Conta sottovoce fino a cinque, solleva le dita dalla tastiera: «Avrei bisogno che per dieci minuti steste in riga, brave e ordinate, ok? Poi rispondo a tutte.»\\
Il sorriso di Valentina illumina la stanza pi\`u del monitor di Laura: «Giochiamo all’intrusione? Ti prego, Tata, fai le vocine.»\\
«Non mi hai capito, Valentina», la voce \`e ferma e leggermente bassa, «ho detto che abbiamo i minuti contati, non abbiamo tempo di giocare.» Guarda una per una le tre ragazze: «Per questo abbiamo bisogno di... pianificare perfettamente l’operazione. Mappare l’infrastruttura per identificare i punti deboli e sondare il terreno con una perlustrazione a bassa intensità.»\\
A quanto pare la missione \`e iniziata e il capitano della squadra \`e indiscutibilmente Laura. I comandi diventano armi, il terminale una mappa, le righe di codice corridoi da attraversare. Ma dietro il gioco la partita \`e vera: non contro qualcuno, ma contro il tempo.
\par\medskip
\begin{center}
  \includegraphics[width=0.7\linewidth]{piantina casa laura senza sfondo.png}
\end{center}

«Sì!» la incita Valentina. Alice si avvicina alla sorella e segue la scena con lei.

\par\medskip
Laura esegue una scansione delle porte del vecchio sistema informatico della polizia locale.

\par\medskip
«Bene. Procediamo con uno \textit{script} veloce. Prendiamo il binocolo tattico e osserviamo le linee nemiche...»

\par\medskip
«Vai Laura!»

\par\medskip
«Fuoco!»

\par\medskip
Laura \`e decisa e schiaccia il tasto INVIO con determinazione ma senza forza. La sensibilità \`e importante. Un comando non deve partire per sbaglio.
\par\medskip
\begin{terminalbox}
laura@dpl:~$ nmap -sS -p21,80 35.228.xxx.xxx
Starting Nmap 7.80 ( https://nmap.org ) at 2041-10-03 23:32 CEST
Nmap scan report for 35.228.xxx.xxx
Host is up (0.088s latency).

PORT   STATE SERVICE
21/tcp open  ftp
80/tcp open  http

laura@dpl:~$
\end{terminalbox}

\par\medskip
«Bene. Abbiamo trovato. Ci sono due porte aperte. Procediamo.»

\par\medskip
Alice sta stringendo la mano a Caterina. Chissà se se ne \`e accorta o se lo ha fatto automaticamente. Ma mentre racconto, Laura \`e gi\`a pronta per la sua prossima mossa.

\par\medskip
«Prepariamo il grimaldello... questo no, troppo invadente... questo no...», Laura verifica diversi comandi al terminale. «Ah! Eccolo, il mio preferito, si insinua come un'\textit{Idra} e  apre le porte ordinarie senza farsi sentire.»\\

\par\medskip
«Cosa fai, Tata?»

\par\medskip
«Abbiamo trovato un varco nella loro recinzione, una porta mal sorvegliata e adesso proviamo a vedere che serratura ha.»

\par\medskip
Il silenzio si potrebbe pesare. Anche Rocky \`e immobile, ma forse gi\`a dorme. Alice fissa Laura che intanto ha ottenuto qualche nuova informazione:\\
\begin{terminalbox}
laura@dpl:~$ hydra -l ftpuser -P password.txt ftp://35.228.xxx.xxx
Hydra v9.2 starting at 2041-10-03 23:34:11

[DATA] max 16 tasks per 1 server
[DATA] 102 login tries: 1 user / 102 passwords
[DATA] attacking ftp://35.228.xxx.xxx:21/

[21][ftp] host: 35.228.xxx.xxx   login: ftpuser   password: ftppass

1 of 1 target successfully completed.
1 valid password found.

Hydra finished at 2041-10-03 23:34:19
laura@dpl:~$
\end{terminalbox}

«Davvero? Queste credenziali cancellano i miei sensi di colpa. La polizia locale se l’\`e cercata.»\\
Caterina sgrana gli occhi e apre leggermente le labbra. Ma Laura le strizza l’occhio e le sussurra:
«\`E un gioco, Cate. Non faremo danni.»

\separatore

\par\medskip
La sala di controllo della polizia  \`e tranquilla. \`E quasi mezzanotte. Gli agenti sono al termine del loro turno. Ore passate a sorvegliare sistemi informatici diventati obsoleti, ma ancora funzionanti. Sottobosco digitale, pullulante di attività tutto sommato innoqua e di scarsa rilevanza.\\
Un agente nota un cursore cambiare colore e spostarsi a destra lasciando una scia di lettere alla sua sinistra. Si avvicina, osserva il monitor:
«Attacco BruteForce sul server 35.228...»\\
Il capitano sta prendendo un caff\`e. Lo appoggia, si avvicina all’agente: «Fa’ vedere. \`E una cosa seria?»\\
Sbadiglia: «Identificate l’IP e vediamo da dove arriva. Sarà ancora uno di quei retrò.»

